% TOP_PROBLEM: {"id":30007666,"problem_number":"LOCAL-30007666","title":"Breaking clientelistic equilibria","category_id":21,"category":{"id":21,"name":"economics","display_name":"Economics","description":"Open economic theory statements, published research agendas, and proposed scoped specifications; question provenance and historical priority are qualified per record.","slug":"economics","order_index":21,"created_at":"2026-10-08T00:00:00Z"},"status":"research_question","external_url":null,"legacy_ids":[],"published":false,"statement":"In the explicitly specified setting: Which reforms durably replace vote buying and broker-controlled benefits with programmatic political competition, and why do clientelistic networks adapt or persist? The mathematical statement above fixes the target, assumptions and scope of this catalogue formulation.","economics":{"original_id":155,"field":"Political economy and institutions","kind":"design","formalization_basis":"adapted_research_question","source_status":"research_agenda","priority_status":"earliest_found","priority_note":"Earliest directly inspected qualifying passage in this audit; earlier origins were not established. A review or research-agenda statement does not imply original priority.","earliest_verified_reference":{"authors":["Jimmy Graham","Horacio Larreguy","Pablo Querubín"],"title":"Clientelism: How It Works, Why It Persists and How to Break It","year":2026,"url":"https://sites.google.com/site/pabloquerubin/research","doi":"10.3386/w34761","locator":"Research page, abstract under the exact chapter title, final sentence","evidence_summary":"The author abstract identifies interventions weakening clientelism and programmatic competition as part of a future research agenda."},"current_reference":{"authors":["Jimmy Graham","Horacio Larreguy","Pablo Querubín"],"title":"Clientelism: How It Works, Why It Persists and How to Break It","year":2026,"url":"https://sites.google.com/site/pabloquerubin/research","doi":"10.3386/w34761","locator":"Research page, abstract under the exact chapter title, final sentence","evidence_summary":"The author abstract identifies interventions weakening clientelism and programmatic competition as part of a future research agenda."},"formalization_note":"The cited passage establishes a research direction, not this exact mathematical statement. The notation, population, policy menu, assumptions, and estimands are an original adaptation for this catalogue; no universal unsolved theorem or source priority is claimed."},"statusQualification":"Published question or research agenda verified at the cited locator. The mathematical specification is adapted for this catalogue and includes additional assumptions; source publication does not establish first-ever priority or certify this exact model remains unsolved. The cited passage establishes a research direction, not this exact mathematical statement. The notation, population, policy menu, assumptions, and estimands are an original adaptation for this catalogue; no universal unsolved theorem or source priority is claimed.","statusReviewedAt":"2026-10-08"}
% FIRST_POSTED: 2026-10-08
% ENTRY_KIND: statement_only
% !TeX program = lualatex
\documentclass[11pt]{article}
\usepackage{fontspec}
\usepackage[a4paper,margin=25mm]{geometry}
\usepackage{amsmath,amssymb,mathtools}
\usepackage{xurl}
\usepackage[hidelinks]{hyperref}
\newcommand{\sref}[2]{\hyperlink{source:#1:#2}{[S#2]}}
\setlength{\emergencystretch}{3em}
\begin{document}
% problemId:30007666
\section*{Breaking clientelistic equilibria}\label{30007666}
\subsection{Definitions and mathematical statement}
Fix a political jurisdiction with politicians, brokers, and voters. Define clientelism as targeted material benefits exchanged for votes or electoral participation under specified monitoring arrangements. Let \(p\) be a feasible intervention changing benefit eligibility, ballot secrecy, broker monitoring, or public information. Let \(C_t(p)\) be the independently measured frequency or value of clientelistic exchange and \(G_t(p)\) public-good provision. Let \(V_t(p)\) measure responsiveness to programmatic performance rather than private transfers, using a predeclared empirical rule. Estimate \(\tau_C(T,p)=E[C_T(p)-C_T(0)]\), together with changes in \(G\) and \(V\), across multiple election cycles. Specify strategic responses by politicians and brokers, financing sources, and substitution into turnout buying, coercion, or patronage. Short-run reductions in reported vote buying do not establish a transition to a durable programmatic equilibrium. Identification requires credible intervention variation, independent behavior measurement, and explicit assumptions about spillovers and secrecy. Determine which bundles alter the incentives sustaining clientelistic exchange and whether effects survive withdrawal. The author's chapter abstract explicitly includes interventions and future research; the particular durability test and outcome vector here are an adaptation rather than an original published conjecture.

\par\textbf{Formulation and scope.} The cited passage establishes a research direction, not this exact mathematical statement. The notation, population, policy menu, assumptions, and estimands are an original adaptation for this catalogue; no universal unsolved theorem or source priority is claimed.

\par\textbf{Open-question provenance.} An explicit question or research agenda was verified at the source locator. This does not by itself certify that every later version remains unresolved \sref{30007666}{1}.

\par\textbf{Historical priority.} Earliest directly inspected qualifying passage in this audit; earlier origins were not established. A review or research-agenda statement does not imply original priority.

\subsection{Short English statement}
In the explicitly specified setting: Which reforms durably replace vote buying and broker-controlled benefits with programmatic political competition, and why do clientelistic networks adapt or persist? The mathematical statement above fixes the target, assumptions and scope of this catalogue formulation.

\subsection{Sources}
\begin{itemize}\raggedright
\item[\textnormal{[S1]}]\hypertarget{source:30007666:1}{} Jimmy Graham, Horacio Larreguy, Pablo Querubín. 2026. Clientelism: How It Works, Why It Persists and How to Break It. Research page, abstract under the exact chapter title, final sentence.
\url{https://sites.google.com/site/pabloquerubin/research}.
DOI: 10.3386/w34761.
{\small\textit{Earliest verified question/agenda reference; historical priority qualified below. The author abstract identifies interventions weakening clientelism and programmatic competition as part of a future research agenda.}}
\end{itemize}
\end{document}
